% !TeX root = ../main.tex
%
% Broader impact statement.  Deliberately short and deliberately unexcited:
% this work makes existing models cheaper to run, and unlocks nothing new.
% Do not inflate it.

\section{Broader impact}
\label{sec:stmt:impact}

Pruning makes an already-trained model cheaper to run. A sparse network needs
less memory and, given hardware or kernels that exploit the sparsity pattern,
less time and less energy per inference than the dense model it came from. The
benefit of that is real but ordinary: lower serving cost, lower emissions per
query, and the possibility of running a model on a device that could not hold
the dense one. This work does not introduce a new capability, a new model or a
new dataset; it studies which training objective produces a better model at a
given sparsity, and the artefacts it produces are ten-class image classifiers.

The risk worth naming is specific to the sparsity regime we study. At 99\% and
99.9\% sparsity, accuracy is lost, and there is no reason for that loss to be
distributed evenly. Capacity removed from a network is removed first from
whatever it represented least robustly, which tends to be the rare cases: rare
classes, rare conditions, and --- in settings other than this benchmark ---
under-represented subgroups. A single aggregate accuracy figure, which is what
this paper and most of the pruning literature report, is exactly the statistic
that cannot see this. We did not measure subgroup effects here; CIFAR-10 is
class-balanced and carries no subgroup annotations, so this benchmark could not
show them even in principle. The practical implication is a caution about
deployment rather than a finding: a model sparsified to these levels for a
consequential application should be evaluated per group and per rare class, not
only in aggregate, and the aggregate figures in this paper should not be read
as evidence that nothing was lost.

We are aware of no path from this work to a directly harmful application that
was not already open. Making inference cheaper marginally lowers the cost of
deploying any classifier, benign or otherwise, but the models studied here are
standard published architectures on a standard public benchmark, and nothing
about the objective is specific to a sensitive use.
